\section{当事人}
\begin{enumerate}
  \item 和案件有直接的利害关系，主动推动诉讼进程的诉讼参与人。
  \item 做题的时候\zd{第一步就要先识别身份}。
  \item 共有的诉讼权利。\zd{提出控告、申请回避、参加法庭调查、使用本民族语言诉讼、申诉}。
  \item \zd{被害人}（\zd{狭义，公诉案件}）。
        \begin{enumerate}
          \item 报案、控告。找公检法均可，检法移交公安。
          \item 公安不立案不服：向原公安申请复议复核 \hz 向同级检察院申诉 \hz 向法院提出自诉。没有先后顺序。
          \item 检察院不起诉不服：向上一级检察院申诉 \hz 向法院提出自诉（公诉转自诉）。
          \item 一审\zd{判决}不服：只能请求检察院抗诉。\zd{被害人是唯一不能上诉的当事人}（公诉是检察院提的）。
          \item 生效裁判不服：再审的申诉。
          \item 委托代理人。自案件\zd{移送审查起诉之日}起。
        \end{enumerate}
  \item 自诉人。在自诉案件中用\zd{自己的名义}，直接向法院提起诉讼的人。大多数案件都是公诉为主，自诉为辅。
        \begin{enumerate}
          \item 直接提起自诉，随时委托诉讼代理人；
          \item 能和解，能撤诉，能调解（公诉转自诉不能），能上诉。
          \item 可以申请法院调取证据。\zd{书面书面，口头（申请）对口头}（回复）。
        \end{enumerate}
  \item 嫌疑人/被告人。提起公诉前/后。
        \begin{enumerate}
          \item 防御性权利。针对\zd{合法}诉讼行为。
                \begin{enumerate}
                  \item 辩护、拒绝回答、开庭前 10 日内收到起诉书副本。
                  \item 参加法庭调查、法庭辩论、\zd{最后陈述权}、\zd{反诉}。
                \end{enumerate}
          \item 救济性权利。针对\zd{不合法、侵权}。对决定、裁决、判决予以审查改变 \hz 撤销。
                \begin{enumerate}
                  \item 申请复议、控告、变更强制措施。
                  \item 申诉（\zd{酌定不起诉决定} \yjt 对检察院；生效裁判 \yjt 对检察院、法院）、上诉、缺席审判异议。
                \end{enumerate}
          \item 其他程序权利。
        \end{enumerate}
  \item 附民原告被告人。类推民诉，\zd{仅限民事部分}行使权利。
  \item 单位当事人。诉讼代表人。被告人地位。
        \begin{enumerate}
          \item 法定代表人 \yjt 未被指控 \hz 不是证人的直接责任人。
          \item 律师担任，只能单位委托。
          \item \zd{不符合规定 \hz 通知 + 不到庭} \yjt 法院要求\zd{检察院另行确定}。
        \end{enumerate}
\end{enumerate}
